\makeabstract
{
Crystallization of Catalase from Bovine Liver
}
{
James Shriver (1), Aidan Cahill (1), Dorothy Huang, PhD (1)
}
{
\textsuperscript{1}Amherst College
}
{
This study attempts to replicate and modify a relatively old method of crystallizing catalase from solid liver. The method that inspired the experiment is from the 1950s and provides a number of ways in which enzyme can be crystallized depending on the type and amount of liver.
}
{
Solid liver was crushed, filtered, and decanted to isolate the enzyme. Two livers were used following the same procedure, although the second was carried out much more efficiently. Acetone was frequently added to the solution in order to precipitate the enzyme. After two rounds of vacuum and gravity filtration, the resulting solution was allowed to dialyze for several days before the enzyme was expected to crystallize. 
}
{
No crystals were detected upon completing our experiment. First, we expected a precipitate to form after several days of dialysis, but this was not the case. There are several reasons that this is possible, with the foremost being that our earlier rounds of filtration were likely not effective enough at removing lipids and other large molecules. Failing to obtain this precipitate, we extracted the solution from the dialysis tubes and attempted to “recrystallize,” a method that would have purified any precipitate that we had collected. This effort also failed to yield crystals. Dr. Kendra Marcus used a microscope to help our search but we were unable to find any crystals. 
}
{
Our experiment failed to isolate and crystallize catalase enzymes after 48 hours, likely due to shortcomings in the early filtration steps. In the future, we may experience better results with more selective filtration as well as by allowing more time for dialysis and recrystallization. 
}

\newpage
